% ============================================================
% III. METHODOLOGY
% ============================================================

\section{Methodology}

The first three exercises were implemented in the \texttt{signal\_tools} package of the \texttt{TEST} project, built on NumPy and SciPy \cite{numpy, scipy}, and executed via \texttt{python main.py}, which generates all figures in this section directly from the project's real data.

\subsection{Item 1 — Analysis of \texttt{Muestra01.csv}}

The file \texttt{Muestra01.csv} contains $1000$ amplitude values, one per line, recorded at a sampling rate $f_s = 5000$ Hz. Following the assignment, the signal was split into segments of $dt = 0.05$ s duration ($250$ samples per segment at 5000 Hz), yielding $4$ complete segments (\texttt{segment\_signal}). For each segment, its spectrum was computed via \texttt{numpy.fft.rfft} and its dominant frequency (\texttt{dominant\_frequency}), excluding the DC component.

\subsection{Item 2 — Sampling and quantization of a 100 Hz sine}

A reference signal $x(t) = \sin(2\pi \cdot 100 \cdot t)$ was generated and sampled (\texttt{sample\_signal}) at three rates: $70$, $500$, and $1000$ Hz, over a $0.05$ s window. Each set of samples was quantized (\texttt{quantize\_dac}) with $b=3$ bits ($8$ levels), simulating a DAC conversion process. For $f_s = 70$ Hz, $f_s < 2\cdot100$ Hz, so the sampling violates the Nyquist theorem and produces aliasing; for $500$ Hz and $1000$ Hz, $f_s \geq 2\cdot100$ Hz and there is no aliasing.

\subsection{Item 3 — FFT of violin, drum, and cat}

Three real WAV recordings were used:

\begin{itemize}
\item \texttt{violin.wav} — $44\,100$ Hz, $2.00$ s
\item \texttt{tambor.wav} — $44\,100$ Hz, $0.96$ s
\item \texttt{gato.wav} — $11\,025$ Hz, $1.58$ s
\end{itemize}

Each file was loaded with \texttt{scipy.io.wavfile} (\texttt{load\_wav}), downmixed to a single channel if stereo, and normalized to $[-1,1]$. The single-sided magnitude spectrum (\texttt{compute\_fft\_spectrum}) and dominant frequency were computed for each full signal.

\subsection{Item 4 — Sampling precision on ESP32-WROOM and Arduino Uno}
\label{sec:metodologia-embebidos}

Per the assignment, this item was designed for two boards: an ESP32-WROOM, which supports all three mechanisms, and an Arduino Uno, which — having a single-core ATmega328P with no FreeRTOS — only supports mechanisms 4(a) and 4(c) (there is no second core to pin a 4(b)-style task to). Each sketch reads three potentiometers at a nominal period of $100$ ms, while the processor concurrently runs a blocking summation from $1$ to $N$ (with $N \in \{1\,000\,000,\ 10\,000\,000,\ 100\,000\,000\}$):

\begin{enumerate}
\item \textbf{4(a) — busy loop}: everything (sampling and summation) happens inside \texttt{loop()}, on a single thread. Implemented for both boards.
\item \textbf{4(b) — dual core}: sampling runs as a FreeRTOS task pinned to core 0 (\texttt{xTaskCreatePinnedToCore}); the summation stays in \texttt{loop()}, running on core 1. ESP32-WROOM only.
\item \textbf{4(c) — interrupt}: a hardware timer triggers an ISR every $100$ ms that only marks a timestamped flag; \texttt{loop()} processes the flag (reads the potentiometers and prints) and then runs the same blocking summation. Implemented for both boards (ESP32's \texttt{timerBegin}/\texttt{timerAttachInterrupt}/\texttt{timerAlarm} API, Arduino-ESP32 core 3.x, vs. the Uno's \texttt{Timer1} in CTC mode).
\end{enumerate}

The three ESP32 sketches were flashed on a real ESP32-WROOM board and their serial output was captured directly with the Arduino IDE's Serial Monitor at $115200$ baud, for each of the three values of $N$. The full captures (eight samples per sketch/$N$ combination) are in \texttt{firmware/real\_logs/*.txt}.

No physical Arduino Uno was available for this project — only the ESP32-WROOM used above, which is a different microcontroller (32-bit Xtensa, dual-core) and cannot run the Uno-specific sketches as written: \texttt{arduino\_uno\_a\_busy\_loop.ino} and \texttt{arduino\_uno\_c\_interrupt.ino} use ATmega328P hardware registers (\texttt{TCCR1A}, \texttt{TCCR1B}, \texttt{OCR1A}, \texttt{TIMSK1}, \texttt{ISR(TIMER1\_COMPA\_vect)}) that do not exist on the ESP32 and would fail to compile for it. The two sketches are written and ready to flash on a real Uno, but the Uno results reported here instead come from \texttt{firmware/simulate\_arduino\_uno\_serial\_output.py}, an explicitly documented timing model (an assumed $\sim\!800\,000$ additions/second on the ATmega328P, a $100.03$ ms \texttt{Timer1} CTC period, $\pm 1$ ms ISR jitter) rather than a hardware capture. The modeled logs, headed \texttt{\# SIMULATED}, are in \texttt{firmware/simulated\_logs\_arduino\_uno/*.txt}.
